\section{Empirical Results}
\label{sec:empirical}

\subsection{Experimental Setup}

We compare three algorithms on three states chosen to span the range of
congressional delegation sizes and geographic complexity:
\begin{itemize}
\item \textbf{NC} (North Carolina, $k=14$, $n \approx 2{,}660$ tracts):
  medium complexity, competitive two-party state.
\item \textbf{WI} (Wisconsin, $k=8$, $n \approx 1{,}900$ tracts):
  small delegation, historically contested.
\item \textbf{TX} (Texas, $k=38$, $n \approx 5{,}300$ tracts):
  large delegation, geographically diverse with major metropolitan areas.
  This is the primary target state where single-chain local-optimum trapping
  is most pronounced.
\end{itemize}

All runs use 2020 Census data, geographic boundary-length edge weights
(B.2 default~\citep{b2paper}), the \geosec{} ratio-optimal tree structure (B.8),
and $T=1{,}000$ steps per chain.
The three algorithms compared are:

\begin{enumerate}
\item \textbf{ForestRecom ($N=1$)}: Single \fr{} chain, $T=1{,}000$ steps,
  $\varepsilon = 0.005$.
  This is the standard baseline.
\item \textbf{PT ($N=2$)}: \pt{} with 2 replicas, cold $\varepsilon=0.005$,
  hot $\varepsilon=0.05$, \texttt{swap\_interval}$=10$.
\item \textbf{PT ($N=4$)}: \pt{} with 4 replicas, cold $\varepsilon=0.005$,
  hot $\varepsilon=0.05$, \texttt{swap\_interval}$=10$.
  Intermediate tolerances: $\varepsilon_1 \approx 0.013$, $\varepsilon_2 \approx 0.033$.
\end{enumerate}

We also report comparison figures for \ms{}~\citep{dellaLibera2026multiscale}
on TX, where available from G.11.

\paragraph{Metrics.}
For each (state, algorithm) pair we report:
\begin{itemize}
\item \textbf{Min EC}: Minimum edge cut observed in the cold chain's record list
  (equivalently, the plan returned at percentile $p=0$).
\item \textbf{Mean EC}: Mean edge cut over all cold chain records.
\item \textbf{Swap Rate}: Observed swap acceptance rate across all adjacent
  replica pairs and all swap opportunities.
\item \textbf{Cold Accept}: Cold chain internal \fr{} acceptance rate.
\item \textbf{Runtime ratio}: Wall-clock time relative to single-chain
  \fr{} ($N=1$), on a single core.
\end{itemize}

\subsection{Results}

Table~\ref{tab:main} presents the comparison.
\pt{}($N=4$) achieves lower minimum edge cut than single-chain \fr{} in all
three states, with the largest improvement on Texas ($k=38$) where local-optimum
trapping is most severe.
The swap acceptance rates of $18$--$24\%$ are consistent with the
${\approx}23\%$ empirical target from \citet{kone2005}, providing evidence that
the geometric ladder is well-calibrated for redistricting chains.

\begin{table}[t]
\centering
\caption{%
  Algorithm comparison on NC ($k=14$), WI ($k=8$), TX ($k=38$).
  Min EC: minimum edge cut over cold chain record list (lower is better).
  Mean EC: mean edge cut over cold chain (lower is better).
  Swap Rate: observed swap acceptance rate.
  Cold Accept: cold chain internal acceptance rate.
  Runtime: wall-clock time relative to single-chain \fr{} ($N=1$).
}
\label{tab:main}
\renewcommand{\arraystretch}{1.25}
\begin{tabular}{@{}llrrrrrr@{}}
\toprule
State & Algorithm & Min EC & Mean EC & Swap Rate & Cold Accept & Runtime \\
\midrule
\multirow{3}{*}{NC ($k=14$)}
 & \fr{} ($N=1$)  & 4{,}821 & 4{,}953 & ---  & 0.47 & $1.0\times$ \\
 & PT ($N=2$)     & 4{,}716 & 4{,}891 & 0.21 & 0.44 & $2.0\times$ \\
 & PT ($N=4$)     & 4{,}663 & 4{,}841 & 0.22 & 0.43 & $4.0\times$ \\
\midrule
\multirow{3}{*}{WI ($k=8$)}
 & \fr{} ($N=1$)  & 2{,}103 & 2{,}174 & ---  & 0.51 & $1.0\times$ \\
 & PT ($N=2$)     & 2{,}059 & 2{,}138 & 0.23 & 0.50 & $2.0\times$ \\
 & PT ($N=4$)     & 2{,}031 & 2{,}110 & 0.23 & 0.49 & $4.0\times$ \\
\midrule
\multirow{3}{*}{TX ($k=38$)}
 & \fr{} ($N=1$)  & 11{,}440 & 11{,}892 & ---  & 0.38 & $1.0\times$ \\
 & PT ($N=2$)     & 10{,}893 & 11{,}421 & 0.19 & 0.36 & $2.0\times$ \\
 & PT ($N=4$)     & 10{,}571 & 11{,}198 & 0.18 & 0.35 & $4.0\times$ \\
\bottomrule
\end{tabular}
\end{table}

\subsection{Analysis of Results}

\paragraph{Large-$k$ benefit.}
The EC improvement from \fr{}($N=1$) to PT($N=4$) is 3.3\% for NC, 3.4\% for
WI, and 7.6\% for TX.
The substantially larger improvement on TX confirms the local-optimum trapping
hypothesis: the single chain on TX gets stuck in a basin with EC $\approx 11{,}440$
from which it cannot escape without the diversity injection provided by hot
chain swaps.
After an accepted swap, the cold chain finds itself in a different plan region
and can explore from there; its internal \fr{} proposals (at $\varepsilon=0.5\%$)
now find lower-EC plans that were inaccessible from the original basin.

\paragraph{Replica count scaling.}
Adding replicas beyond $N=2$ provides diminishing but positive returns.
PT($N=2$) already captures most of the benefit: the step from $N=1$ to $N=2$
accounts for approximately $60$--$70\%$ of the total PT($N=4$) improvement.
The step from $N=2$ to $N=4$ adds a further 30--40\%.
This is consistent with the empirical observation that a two-replica system
already achieves substantial mixing improvement when the ladder is geometric;
adding replicas further refines the coverage but with diminishing marginal
returns~\citep{earl2005}.

\paragraph{Swap acceptance rates.}
The observed swap rates of $18$--$24\%$ are near the ${\approx}23\%$ empirical
target.
The TX chains show slightly lower swap rates ($18$--$19\%$) because the larger
state has more edge-cut variance across replicas---each replica's current EC is
drawn from a wider distribution, which reduces the probability that a cold-hot
swap is accepted.
This is consistent with the Kone-Kofke theory: wider EC distributions between
replicas correspond to a less ideal geometric spacing.
For TX with $N=6$ replicas (as recommended in the parameter scaling table of the
spec), the swap rates are expected to return to the $21$--$24\%$ range.

\paragraph{Cold chain acceptance rates.}
The cold chain acceptance rates ($0.35$--$0.51$) are healthy across all states.
The lowest rate is TX ($0.35$ for PT($N=4$)): the cold chain at $\varepsilon=0.5\%$
has a somewhat lower \fr{} acceptance rate on Texas than on NC or WI, because
the balance constraint is harder to satisfy in a large state with concentrated
metropolitan populations.
This rate is not pathologically low; a rate below $0.10$ would indicate that
the cold tolerance is too tight for the state's population distribution.

\paragraph{Comparison with Multi-Scale.}
For TX, \ms{}~\citep{dellaLibera2026multiscale} achieves a minimum EC of
approximately $10{,}380$ (from G.11), compared to PT($N=4$)'s $10{,}571$.
\ms{} has an edge on TX because the geographic coarsening allows larger structural
changes per step---it effectively operates at a coarser spatial resolution
that spans the metropolitan--suburban boundary, which is the primary
local-optimum barrier on Texas maps.
PT is competitive with \ms{} and does not require a coarsening heuristic;
for states where geographic coarsening is difficult to calibrate (complex
coastlines, archipelago counties), PT may be preferred.

\paragraph{NC and WI: modest but positive improvement.}
For NC ($k=14$) and WI ($k=8$), the single-chain \fr{} already mixes
adequately and PT provides a modest but positive improvement of 3--4\%.
This is expected: when the chain mixes well without hot chain assistance, the
replica exchange mechanism still occasionally injects useful diversity from
the hot chain's broader sampling of plan space.
The cost is $N\times$ the runtime; for NC and WI, the runtime ratio of $4\times$
for PT($N=4$) is a significant overhead for a modest gain.
For these states, single-chain \fr{} or \cs{} (convergence sweep) is the
recommended search strategy.
